% Generated draft: requires mathematical and source review.
\documentclass[11pt]{article}
\usepackage[margin=25mm]{geometry}
\usepackage{fontspec}
\setmainfont{Latin Modern Roman}
\usepackage{amsmath,amssymb,mathtools,unicode-math}
\setmathfont{Latin Modern Math}
\usepackage{enumitem,needspace,xurl,hyperref}
\setlength{\parindent}{0pt}
\setlength{\parskip}{5pt}
\setlength{\emergencystretch}{3em}
\newcommand{\N}{\mathbb{N}}
\newcommand{\Z}{\mathbb{Z}}
\newcommand{\Q}{\mathbb{Q}}
\newcommand{\R}{\mathbb{R}}
\newcommand{\C}{\mathbb{C}}
\newcommand{\F}{\mathbb{F}}
\newcommand{\A}{\mathbb{A}}
\newcommand{\PP}{\mathbb{P}}
\newcommand{\E}{\mathbb{E}}
\newcommand{\eps}{\varepsilon}
\newcommand{\dd}{\,\mathrm{d}}
\begin{document}
\section*{UnsolvedMath ID: 30002050}
\textbf{Source problem code:} catalogue entry
\subsection{Definitions and mathematical statement}
The supplied record retains a conditional theorem schema, not a separately verified open problem. Its verification note marks the statement as unrecoverable as an independent problem: the official contribution is described as asserting the decomposition conditionally on Zilber's conjecture, with no separate clean statement promoted. The formulation below preserves that conditional character.

For a field \(F\), the algebraic torus \(\mathbb G_m^n\) has \(F\)-points \((F^\times)^n\). A Laurent polynomial over \(F\) is a finite sum
\[
f=\sum_{a\in\Z^n}c_a x_1^{a_1}\cdots x_n^{a_n},\qquad c_a\in F,
\]
and its support is \(\operatorname{supp}(f)=\{a\in\Z^n:c_a\ne0\}\). To express geometric zero sets, choose a common algebraically closed field \(\Omega\) containing the input and output coefficients and put
\[
Z_\Omega(f_1,\ldots,f_r)=\{x\in(\Omega^\times)^n:f_i(x)=0\text{ for every }1\le i\le r\}.
\]
This is notation for a geometric reading; the source does not specify the intended field of points, the output coefficient field, or the field over which irreducible components are taken.

\textbf{Supplied conditional decomposition schema.} Let \(K\) be a number field and let \(V\subseteq\mathbb G_m^n\) be a subvariety over \(K\), defined by a bounded number of equations of bounded height, degree at most \(d\), and bounded monomial support, with these bounds understood in the unspecified sense intended by the source. Assuming the intended form of Zilber's conjecture, there is an algorithm producing a finite indexing set \(A\) and, for each \(\alpha\in A\), a sequence of Laurent polynomials
\[
(P_{\alpha,1},\ldots,P_{\alpha,L_\alpha},Q_\alpha),\qquad L_\alpha\le n,
\]
such that, in the geometric interpretation just described,
\[
V_\Omega=\bigcup_{\alpha\in A}U_\alpha,\qquad
U_\alpha=Z_\Omega(P_{\alpha,1},\ldots,P_{\alpha,L_\alpha})\setminus Z_\Omega(Q_\alpha).
\]
Every irreducible component \(Y\) of every nonempty \(U_\alpha\) has ambient codimension \(L_\alpha\); equivalently, with closure taken in \(\mathbb G_{m,\Omega}^n\),
\[
n-\dim\overline{Y}=L_\alpha.
\]
The stated operation bound is \(O(\log d)\) arithmetic operations, interpreted asymptotically as \(d\) grows.

\textbf{Unresolved qualifications.} The original statement says only ``Zilber's conjecture,'' whereas the supplied literature assessment describes the premise as an effective Zilber conjecture. Neither version is stated. The meanings of the height and support bounds, the degree convention, the parameters held fixed in the operation bound, and the computational and output representations are not supplied. The field conventions above and the possibility of an empty list of \(P\)'s also require confirmation. Consequently this is a conditional schema with explicit gaps, not a fully specified algorithmic theorem or an unconditional conjecture.
\subsection{Short English statement}
The retained claim says that, assuming the intended Zilber conjecture, a sparsely described subvariety of a torus can be decomposed into finitely many locally closed pieces, each defined by as many equations as its codimension, using \(O(\log d)\) arithmetic operations. Essential hypotheses and complexity conventions remain unspecified. The supplied verification note treats this as a conditional theorem rather than an additional open problem.
\subsection{Sources}
\begin{itemize}
\item Local UnsolvedMath registry, permanent ID \(30002050\), exact source problem code \texttt{catalogue entry}, title ``Conditional Sparse-Decomposition Theorem,'' in the Oberwolfach Reports Open Problems collection. No individual proposer is supplied; the proposed year is recorded as 2012.
\item Supplied report citation: \emph{Diophantische Approximationen (2013)}, Oberwolfach Reports, DOI 10.4171/owr/2012/22, pp.\ 1307--1359. \url{https://doi.org/10.4171/owr/2012/22}. The background describes these reports as documenting mathematical workshops and their open-problem discussions. The citation's 2013 label and the recorded 2012 provenance require reconciliation.
\item Amoroso--Leroux--Sombra, \emph{Overdetermined systems of sparse polynomial equations}. Supplied literature source: \url{https://arxiv.org/abs/1307.5788}.
\item Capuano, \emph{Exploring Zilber's conjecture on atypical intersections}. Supplied literature source: \url{https://doi.org/10.1112/jlms.70536}.
\item Supplied literature assessment checked on \texttt{2026-08-21}: the registry labels the record open and states that the displayed theorem depends on an effective Zilber conjecture whose general premise remains unresolved, so the extracted unconditional claim is not settled. This is supplied dated curation derived from the OWR material, not an independently verified literature assessment. The separate supplied verification note says that the conditional theorem is not an additional open problem beyond the conjecture retained in a preceding record. No source was independently consulted for this draft.
\end{itemize}
\end{document}
